% Section 1 -- Introduction (rewritten 2026-11-23, condensed the same day). Narrative modelled on DeGrave et al. 2021
% (promise -> data collection creates worst-case confounding -> consequence -> what we show) and Song et al. 2020.
\section{Introduction}\label{sec:intro}

% P1 background
Two neurodevelopmental disorders, autism spectrum disorder (ASD) and attention deficit hyperactivity disorder (ADHD), are diagnosed from behaviour, through interviews and observation that are slow and depend on the examiner. Many studies therefore train machine learning models on functional connectivity to find objective markers, and because a single hospital rarely has enough participants, they pool data across sites, as in the ABIDE and ADHD-200 consortia~\cite{dimartino2014,bellec2017}.

% P2 the general risk of pooling
Pooling has a price. Sites differ in scanners, protocols and recruitment, and these differences can be as large as the effects of the disorders~\cite{yamashita2019}. They are also easy to detect: a classifier identified which of $17$ data sets a brain scan came from with an accuracy above $70\%$, where chance was $6\%$~\cite{wachinger2021}. If the site is related to the diagnosis, a model can use the site as a cue and fail at other sites, known as shortcut learning~\cite{geirhos2020,zech2018}.

% P3 the specific problem
This paper studies the most extreme form of this risk. Suppose that a child psychiatry clinic scans children with autism, while typically developing children are scanned at a university laboratory in another city. Each site then contributes participants of a single diagnosis, so knowing the site means knowing the label; we call such sites \emph{class-pure}. They arise naturally, because patients are recruited where they are treated and controls where they are easiest to reach, and they occur in public consortia: one ABIDE-II collection contains only autistic participants, and one ADHD-200 site only controls. Early COVID-19 chest radiograph collections had the same structure, and models trained on them relied on cues to the data source~\cite{degrave2021}.

% P4 why it matters
Subject-level cross-validation cannot reveal this problem: a held-out patient comes from a site that contributed only patients to training, so recognising the site is enough to classify the patient correctly. In our ABIDE-I experiments, subject-level cross-validation reported an area under the ROC curve (AUC) of $0.94$ for a model that reached only $0.58$ at new sites. Accuracies reported from class-pure data may thus owe a substantial part to the site.

% P5 the gap
Harmonization methods such as ComBat protect the diagnosis as a covariate~\cite{johnson2007,fortin2018}, but at a class-pure site diagnosis and site cannot be told apart, so this protection cannot be applied. Debiasing methods in pattern recognition reweight groups or learn from \emph{bias-conflicting} samples~\cite{sagawa2020,zhang2022cnc,wang2025mdn}; here these would be patients at a control site, and there are none. Exact site invariance~\cite{ganin2016,dinsdale2021,kunda2023} would remove the information that separates the classes~\cite{zhao2019}. In genomics, designs in which batch and group coincide are known to defeat both batch correction and prediction across batches~\cite{song2020,parker2012}, but that work concerns laboratory batches. For diagnostic models trained on many class-pure sites, it remains open how much class purity costs, how such a model should be evaluated, and which data to add and how to train with them.

% P6 what we do
Our method acts on the training data, the training objective and the evaluation, so it can be used with standard score-based classifiers. When every site holds one diagnosis, every comparison between a patient and a control is also a comparison between two sites, and only comparisons within a site are free of site effects. We formalize this with an additive model of site effects and derive three rules (Fig.~\ref{fig:overview}): \emph{evaluate} by holding out whole sites and testing whether the model ranks by site; \emph{design} the data by recruiting a few \emph{bridging} subjects of the missing diagnosis at each site, guided by a coefficient computed from the table of subjects per site and diagnosis; and \emph{estimate} the model only from comparisons within sites. To measure these effects against a known truth, we build class-pure training sets from sites that hold both diagnoses and use the unused participants of the other diagnosis to test whether the model has learned the site. We tested each rule on resting-state fMRI from two disorders in three cohorts, with criteria written down before the corresponding experiments.

% P7 significance
From a pattern recognition perspective, class-pure sites are a spurious correlation without any bias-conflicting samples, and our results show that collecting a few such samples by design can break it.

Our contributions are as follows.
\begin{itemize}
\item \textbf{Analysis.} With an additive site-effect model, we show that class-pure data cannot separate the diagnosis from the site without an untestable assumption, and that adding subjects to the same sites does not reduce the site part of the error. A design coefficient $L$, computed from the site-by-class table, measures how much of it a design retains.
\item \textbf{Evaluation.} Subject-level cross-validation overstated class-pure logistic models by $0.37$ to $0.46$ in AUC across three cohorts, and two published multi-site methods failed in the same setting. Holding out whole sites removed the systematic bias, and a \emph{reverse-set} test on held-out subjects of the opposite diagnosis at the training sites revealed shortcuts that other checks missed.
\item \textbf{Design.} $L$ is a data acquisition criterion that can be computed before any data are collected: it ranked designs by the size of the shortcut, with rank correlations of $-0.970$ to $-0.997$ with the reverse-set AUC. A few bridging subjects of the other class per site curbed the shortcut more than the same number of same-class additions (reverse-set AUC higher by $0.10$ to $0.29$).
\item \textbf{Estimation.} We propose a logistic loss on within-site patient and control pairs, a learning objective whose solution is unchanged by additive site shifts (Proposition~\ref{prop:within}). It kept the reverse-set AUC between $0.59$ and $0.67$ in all $12$ bridged cells of ABIDE-I, and with two-way bridging its test AUC was within $0.01$ of the pooled model; as a linear head on a frozen self-supervised encoder it reached a reverse-set AUC of $0.59$, against at most $0.39$ for deep baselines. On ABIDE-II all assessable pre-specified criteria were met, and on ADHD-200 all assessable criteria except the cost of class purity.
\end{itemize}
